\documentclass[reqno,11pt]{amsart}
\usepackage{math327-lecture}
\usepackage{needspace}
\lectureinfo{9}{Sep. 23, 2026}{Cosets and Partitions}

% Based on handwritten Lecture 7, Sep. 17, 2025.
\begin{document}
\makelecturetitle

\section{Review: normal subgroups}

Last time we introduced normal subgroups. Recall that $N\leq G$ is
\emph{normal} if
\[
  gNg^{-1}\subseteq N\qquad\text{for every }g\in G.
\]
We write $N\trianglelefteq G$.

\begin{proposition}
If $N\trianglelefteq G$, then $gNg^{-1}=N$ for every $g\in G$.
\end{proposition}
\begin{proof}
The definition gives one inclusion. For the other, fix $n\in N$ and put
$n'=g^{-1}ng$. Normality applied to $g^{-1}$ gives $n'\in N$. Since
$n=gn'g^{-1}$, we have $n\in gNg^{-1}$. Thus $N\subseteq gNg^{-1}$.
\end{proof}

\begin{example}
Recall
\[
  S_3=\{1,x,x^2,y,xy,x^2y\},\qquad x^3=1,\quad y^2=1,\quad yx=x^2y.
\]
For $N=\{1,x,x^2\}$, the calculations from Lecture 8 give
\[
  yNy^{-1}=\{1,yxy^{-1},yx^2y^{-1}\}=\{1,x^2,x\}=N.
\]
Multiplying both sets on the right by $y$ gives $yN=Ny$.
\end{example}

More generally, for any subgroup $N\leq G$,
\begin{equation}
  gNg^{-1}=N\quad\Longleftrightarrow\quad gN=Ng.
\end{equation}
This suggests studying the sets $gH$ and $Hg$ for an arbitrary subgroup
$H\leq G$, whether or not $H$ is normal.

\section{Left and right cosets}

\begin{define}[Cosets and representatives]
Let $H\leq G$ and $g\in G$. The \emph{left coset} and \emph{right coset}
of $H$ represented by $g$ are
\begin{equation}
  gH=\{gh:h\in H\},\qquad Hg=\{hg:h\in H\}.
\end{equation}
We call $g$ a \emph{coset representative}. The set of all left cosets is
\begin{equation}
  G/H=\{gH:g\in G\}.
\end{equation}
\end{define}

A coset is a subset of $G$; its representative is an element of $G$.
Different representatives can give the same coset. If $g\in H$, then
$gH=H=Hg$, as we will prove below.

\begin{example}[Cosets of the subgroup of order three in $S_3$]
Take $H=\{1,x,x^2\}$. For its elements, we obtain
\begin{align*}
  1H&=\{1,x,x^2\}=H,\\
  xH&=\{x,x^2,x^3\}=H,\\
  x^2H&=\{x^2,x^3,x^4\}=H.
\end{align*}
For the other three elements, use $yx=x^2y$, $yx^2=xy$, and $xyx=y$:
\begin{align*}
  yH&=\{y,yx,yx^2\}=\{y,x^2y,xy\},\\
  xyH&=\{xy,xyx,xyx^2\}=\{xy,y,x^2y\}=yH,\\
  x^2yH&=\{x^2y,x^2yx,x^2yx^2\}=\{x^2y,xy,y\}=yH.
\end{align*}
There are exactly two distinct left cosets:
\[
  S_3/H=\{H,yH\},\qquad S_3=H\sqcup yH.
\]
The symbol $\sqcup$ denotes a disjoint union. Since $H$ is normal, the
right coset represented by any $g$ equals the corresponding left coset.
\end{example}

\begin{remark}
At this stage $G/H$ is a \emph{set} of cosets. It is not a set of chosen
representatives. When $H$ is normal, we will later give $G/H$ a group
operation.
\end{remark}

\section{Cosets form a partition}

\begin{define}[Partition]
A \emph{partition} of a set $X$ is a collection of nonempty subsets
$\{X_\alpha\}_{\alpha\in I}$ such that
\[
  X=\bigcup_{\alpha\in I}X_\alpha,
  \qquad X_\alpha\cap X_\beta=\varnothing\quad(\alpha\neq\beta).
\]
We write $X=\bigsqcup_{\alpha\in I}X_\alpha$.
\end{define}

\begin{lemma}[Equality of cosets]\label{lem:coset-equality}
Let $H\leq G$ and $g_1,g_2\in G$. The following are equivalent:
\begin{enumerate}[label=\textbf{(\arabic*)}]
  \item $g_1H\cap g_2H\neq\varnothing$;
  \item $g_1H=g_2H$;
  \item $g_2^{-1}g_1\in H$;
  \item $g_1^{-1}g_2\in H$.
\end{enumerate}
\end{lemma}
\begin{proof}
First consider the special case $g_2=e_G$. We show that
\[
  gH\cap H\neq\varnothing
  \quad\Longleftrightarrow\quad g\in H
  \quad\Longleftrightarrow\quad gH=H.
\]
If $g\in H$, closure gives $gH\subseteq H$. Conversely, every $h\in H$
can be written $h=g(g^{-1}h)$, with $g^{-1}h\in H$, so $H\subseteq gH$.
Thus $gH=H$, and their intersection is nonempty.
If $z\in gH\cap H$, write $z=gh$ with $h\in H$. Then
$g=zh^{-1}\in H$, which completes the special case.

For general $g_1,g_2$, left multiplication by $g_2^{-1}$ is a bijection
of $G$. Hence
\begin{align*}
  g_1H\cap g_2H\neq\varnothing
  &\quad\Longleftrightarrow\quad
  (g_2^{-1}g_1)H\cap H\neq\varnothing\\
  &\quad\Longleftrightarrow\quad g_2^{-1}g_1\in H\\
  &\quad\Longleftrightarrow\quad (g_2^{-1}g_1)H=H\\
  &\quad\Longleftrightarrow\quad g_1H=g_2H.
\end{align*}
Finally, $g_2^{-1}g_1\in H$ if and only if its inverse
$g_1^{-1}g_2$ belongs to $H$.
\end{proof}

\begin{proposition}
The distinct left cosets of $H$ form a partition of $G$.
\end{proposition}
\begin{proof}
Every $g\in G$ belongs to $gH$, since $g=ge_G$ and $e_G\in H$.
Thus the cosets are nonempty and cover $G$. By
Lemma~\ref{lem:coset-equality}, two cosets that intersect are equal, so
distinct cosets are disjoint.
\end{proof}

If $R\subseteq G$ contains exactly one representative of each left coset,
we may write
\begin{equation}
  G=\bigsqcup_{r\in R}rH.
\end{equation}
The representative set $R$ is not unique. For example, in the preceding
$S_3$ example, both $\{1,y\}$ and $\{x,xy\}$ are representative sets.

\section{Example: cosets of \texorpdfstring{$8\Z$}{8Z}}

In the additive group $G=\Z$, take $H=8\Z$. Cosets are written additively:
\begin{equation}
  a+8\Z=\{a+8k:k\in\Z\}.
\end{equation}
For instance,
\begin{align*}
  0+8\Z&=\{\ldots,-16,-8,0,8,16,\ldots\},\\
  1+8\Z&=\{\ldots,-15,-7,1,9,17,\ldots\},\\
  7+8\Z&=\{\ldots,-9,-1,7,15,23,\ldots\}.
\end{align*}
Arranging the integers by their remainders gives eight disjoint columns:
\[
\begin{array}{c|rrrrrrrr}
  \text{remainder}&0&1&2&3&4&5&6&7\\ \hline
  &\vdots&\vdots&\vdots&\vdots&\vdots&\vdots&\vdots&\vdots\\
  &16&17&18&19&20&21&22&23\\
  &8&9&10&11&12&13&14&15\\
  &0&1&2&3&4&5&6&7\\
  &-8&-7&-6&-5&-4&-3&-2&-1\\
  &-16&-15&-14&-13&-12&-11&-10&-9\\
  &\vdots&\vdots&\vdots&\vdots&\vdots&\vdots&\vdots&\vdots
\end{array}
\]
Each column is one coset. The equality criterion becomes
\begin{equation}
  a+8\Z=b+8\Z
  \quad\Longleftrightarrow\quad a-b\in8\Z
  \quad\Longleftrightarrow\quad a\equiv b\pmod8.
\end{equation}
For example, $0+8\Z=8+8\Z$ and $1+8\Z=17+8\Z$.
Every integer has exactly one remainder in $\{0,1,\ldots,7\}$, so
\[
  \Z=\bigsqcup_{r=0}^7(r+8\Z),\qquad
  \Z/8\Z=\{0+8\Z,1+8\Z,\ldots,7+8\Z\}.
\]
Writing $\overline a=a+8\Z$, we may abbreviate this as
\[
  \Z/8\Z=\{\overline0,\overline1,\ldots,\overline7\}.
\]
Thus an infinite group can have a finite set of cosets.

\end{document}
